\section{Theoretical Analysis of \name{}}
\label{app:dfb-theory}

This appendix derives the \name{} representation objective and characterizes its population fixed point. The analysis concerns policy evaluation for a fixed task-conditioned policy $\pi_z$. In particular, the contraction result below applies to the successor-density update with $\pi_z$ held fixed; it does not imply convergence of the coupled representation and actor updates used in practice.

Throughout, $\rho$ denotes the dataset next-state marginal. We assume that, for the state-action pairs and policies of interest, $P(\cdot\mid s,a)\ll\rho$ and $\nu^{\pi_z}(s,a,\cdot)\ll\rho$. For the contraction argument, we additionally assume that the relevant density ratios are bounded. We also assume that the partition functions appearing below are finite and nonzero.


\subsection{Normalized successor density}
\label{app:normalized-density}

For a fixed policy $\pi_z$, recall the discounted successor measure
\begin{equation}
    M^{\pi_z}(s,a,ds^+) := \E_{\pi_z}\left[\sum_{t=0}^{\infty}\gamma^t\delta_{s_{t+1}}(ds^+) \,\middle|\, s_0=s,\ a_0=a\right].
    \label{eq:app-sm}
\end{equation}
Its total mass is $1/(1-\gamma)$. \name{} works with the normalized successor measure
\begin{equation}
    \nu^{\pi_z}(s,a,ds^+) := (1-\gamma)M^{\pi_z}(s,a,ds^+),
    \label{eq:app-normalized-sm}
\end{equation}
which is a probability measure over future states.

Multiplying the Bellman recursion for $M^{\pi_z}$ by $1-\gamma$ gives
\begin{equation}
    \nu^{\pi_z}(s,a,ds^+) = (1-\gamma)P(ds^+\mid s,a) + \gamma\E_{\substack{s'\sim P(\cdot\mid s,a)\\a'\sim\pi_z(\cdot\mid s')}}\left[\nu^{\pi_z}(s',a',ds^+)\right].
    \label{eq:app-normalized-bellman}
\end{equation}

Define the one-step transition density ratio
\begin{equation}
    \chi(s,a,s^+) := \frac{dP(\cdot\mid s,a)}{d\rho}(s^+),
    \label{eq:app-transition-ratio}
\end{equation}
and the normalized successor-density ratio
\begin{equation}
    h_z^\pi(s,a,s^+) := \frac{d\nu^{\pi_z}(s,a,\cdot)}{d\rho}(s^+).
    \label{eq:app-true-ratio}
\end{equation}
Because $\nu^{\pi_z}$ is normalized,
\begin{equation}
    \E_{s^+\sim\rho}\left[h_z^\pi(s,a,s^+)\right] = 1.
    \label{eq:app-ratio-normalization}
\end{equation}

Dividing Equation~\eqref{eq:app-normalized-bellman} by $\rho$ yields the Bellman equation for the density ratio:
\begin{equation}
    h_z^\pi(s,a,s^+) = (1-\gamma)\chi(s,a,s^+) + \gamma\E_{\substack{s'\sim P(\cdot\mid s,a)\\a'\sim\pi_z(\cdot\mid s')}}\left[h_z^\pi(s',a',s^+)\right].
    \label{eq:app-density-bellman}
\end{equation}
Thus \name{} can be viewed as policy evaluation in the space of normalized successor-density ratios.


\subsection{Conditional density-ratio estimation}
\label{app:density-estimation}

We next show why the normalized exponential model used by \name{} recovers a density ratio.

Consider any probability measure $\mu\ll\rho$, and let
\[
    g(s^+) := \frac{d\mu}{d\rho}(s^+).
\]
For a real-valued score $f$, define the normalized model
\begin{equation}
    h_f(s^+) := \frac{\exp f(s^+)}{\E_{\tilde s^+\sim\rho}\left[\exp f(\tilde s^+)\right]}.
    \label{eq:app-generic-density}
\end{equation}

\paragraph{Lemma 1 (Density-ratio identification).}
For fixed $\mu$ and $\rho$, consider
\begin{equation}
    \mathcal J_{\mathrm{ratio}}(f) := \E_{s^+\sim\mu}[f(s^+)] - \log\E_{s^+\sim\rho}\left[\exp f(s^+)\right].
    \label{eq:app-ratio-objective}
\end{equation}
Then
\begin{equation}
    \mathcal J_{\mathrm{ratio}}(f) = C_\mu - \operatorname{KL}\left(\mu\,\middle\|\,h_f\rho\right),
    \label{eq:app-ratio-kl}
\end{equation}
where $C_\mu=\E_{\mu}[\log g]$ does not depend on $f$. Hence any maximizer satisfies
\begin{equation}
    h_f(s^+) = \frac{d\mu}{d\rho}(s^+)
    \label{eq:app-ratio-optimum}
\end{equation}
almost everywhere under $\rho$.

\paragraph{Proof.}
Since
\[
    \log h_f(s^+) = f(s^+) - \log\E_{\rho}\left[e^f\right],
\]
we have
\begin{align}
    \operatorname{KL}\left(\mu\,\middle\|\,h_f\rho\right)
    &= \E_{\mu}\left[\log\frac{g(s^+)}{h_f(s^+)}\right] \nonumber\\
    &= \E_{\mu}[\log g(s^+)] - \E_{\mu}[f(s^+)] + \log\E_{\rho}[e^f].
    \label{eq:app-kl-expansion}
\end{align}
Rearranging gives Equation~\eqref{eq:app-ratio-kl}. Since KL divergence is minimized at zero, the optimum satisfies $h_f=g$.
\hfill$\square$

When a finite optimal score exists, Lemma~1 implies
\begin{equation}
    f^*(s^+) = \log\frac{d\mu}{d\rho}(s^+) + c,
    \label{eq:app-optimal-score}
\end{equation}
where the additive constant $c$ disappears under normalization.

\name{} applies this construction conditionally on $(s,a,z)$. We use the bilinear score
\begin{equation}
    f_\theta(s,a,z,s^+) := \frac{F_\theta(s,a,z)^\top B_\theta(s^+)}{\tau},
    \label{eq:app-bilinear-score}
\end{equation}
which induces
\begin{equation}
    h_\theta(s,a,z,s^+) := \frac{\exp f_\theta(s,a,z,s^+)}{\E_{\tilde s^+\sim\rho}\left[\exp f_\theta(s,a,z,\tilde s^+)\right]}.
    \label{eq:app-density-model}
\end{equation}
Thus $F_\theta^\top B_\theta/\tau$ parameterizes a log density ratio, while its normalized exponential represents the successor-density ratio itself.


\subsection{\name{} as a Bellman operator}
\label{app:bellman-operator}

For fixed $z$ and policy $\pi_z$, define the operator
\begin{equation}
    (\mathcal T_z h)(s,a,s^+) := (1-\gamma)\chi(s,a,s^+) + \gamma\E_{\substack{s'\sim P(\cdot\mid s,a)\\a'\sim\pi_z(\cdot\mid s')}}\left[h(s',a',s^+)\right].
    \label{eq:app-dfb-operator}
\end{equation}

Equation~\eqref{eq:app-density-bellman} immediately shows that
\begin{equation}
    h_z^\pi = \mathcal T_z h_z^\pi.
    \label{eq:app-fixed-point}
\end{equation}

The operator also preserves normalization. If $\E_\rho[h(s,a,\cdot)]=1$ for every $(s,a)$, then
\begin{align}
    \E_{s^+\sim\rho}\left[(\mathcal T_z h)(s,a,s^+)\right]
    &= (1-\gamma)\E_\rho[\chi(s,a,s^+)] + \gamma\E_{s',a'}\E_\rho[h(s',a',s^+)] \nonumber\\
    &= (1-\gamma)+\gamma = 1.
    \label{eq:app-operator-normalization}
\end{align}

\paragraph{Proposition 1 (Contraction).}
On the space of bounded density-ratio functions equipped with the supremum norm, $\mathcal T_z$ is a $\gamma$-contraction:
\begin{equation}
    \left\|\mathcal T_z h_1-\mathcal T_z h_2\right\|_\infty \leq \gamma\left\|h_1-h_2\right\|_\infty.
    \label{eq:app-contraction}
\end{equation}

\paragraph{Proof.}
The one-step term cancels, giving
\begin{align}
    \left|(\mathcal T_z h_1)(s,a,s^+) - (\mathcal T_z h_2)(s,a,s^+)\right|
    &= \gamma\left|\E_{s',a'}\left[h_1(s',a',s^+)-h_2(s',a',s^+)\right]\right| \nonumber\\
    &\leq \gamma\E_{s',a'}\left[\left|h_1(s',a',s^+)-h_2(s',a',s^+)\right|\right] \nonumber\\
    &\leq \gamma\left\|h_1-h_2\right\|_\infty.
    \label{eq:app-contraction-proof}
\end{align}
Taking the supremum over $(s,a,s^+)$ yields Equation~\eqref{eq:app-contraction}.
\hfill$\square$

\paragraph{Corollary 1 (Unique fixed point).}
Under the assumptions above, $\mathcal T_z$ has a unique bounded fixed point, equal to the normalized successor-density ratio $h_z^\pi$. Consequently, exact repeated Bellman updates
\begin{equation}
    h_{k+1} = \mathcal T_z h_k
    \label{eq:app-exact-iteration}
\end{equation}
converge to $h_z^\pi$ for any bounded initialization.

This result characterizes the representation update as policy evaluation: for a fixed $\pi_z$, the desired solution is the normalized discounted future-state distribution induced by that policy.


\subsection{Population \name{} objective}
\label{app:population-objective}

We now connect the Bellman operator to the objective optimized by \name{}. Let $\bar h$ be a fixed target density ratio satisfying $\E_\rho[\bar h(s,a,z,\cdot)]=1$, and define
\begin{equation}
    h_z^{\mathrm{tgt}}(s,a,s^+) := (1-\gamma)\chi(s,a,s^+) + \gamma\E_{\substack{s'\sim P(\cdot\mid s,a)\\a'\sim\pi_z(\cdot\mid s')}}\left[\bar h(s',a',z,s^+)\right].
    \label{eq:app-target-density}
\end{equation}
By Equation~\eqref{eq:app-operator-normalization}, $h_z^{\mathrm{tgt}}$ is itself a normalized density ratio.

The corresponding population cross-entropy loss is
\begin{equation}
    \mathcal L_{\mathrm{DFB}}^{\mathrm{pop}}(\theta;\bar h) := -\E_{\substack{z\sim p_Z\\(s,a)\sim\mathcal D_{SA}\\s^+\sim\rho}}\left[h_z^{\mathrm{tgt}}(s,a,s^+)\log h_\theta(s,a,z,s^+)\right].
    \label{eq:app-population-ce}
\end{equation}

\paragraph{Proposition 2 (Bellman update by cross-entropy minimization).}
If the model class can represent $h_z^{\mathrm{tgt}}$ exactly, a global minimizer of Equation~\eqref{eq:app-population-ce} satisfies
\begin{equation}
    h_\theta(s,a,z,s^+) = h_z^{\mathrm{tgt}}(s,a,s^+) = (\mathcal T_z\bar h)(s,a,s^+)
    \label{eq:app-ce-bellman}
\end{equation}
almost everywhere under $\mathcal D_{SA}\times\rho$.

\paragraph{Proof.}
For each fixed $(s,a,z)$, Equation~\eqref{eq:app-population-ce} is the cross-entropy between the target distribution $h_z^{\mathrm{tgt}}\rho$ and model distribution $h_\theta\rho$. It can be written as
\[
    H(h_z^{\mathrm{tgt}}\rho) + \operatorname{KL}\left(h_z^{\mathrm{tgt}}\rho\,\middle\|\,h_\theta\rho\right),
\]
where the first term is independent of $\theta$. The minimum is therefore attained when $h_\theta=h_z^{\mathrm{tgt}}$.
\hfill$\square$

Combining Proposition~1 and Proposition~2, exact optimization against the current density target implements one step of a $\gamma$-contractive policy-evaluation operator.

To recover the score-space objective used in the derivation, substitute
\[
    \log h_\theta(s,a,z,s^+) = f_\theta(s,a,z,s^+) - \log Z_\theta(s,a,z),
\]
where
\begin{equation}
    Z_\theta(s,a,z) := \E_{s^+\sim\rho}\left[\exp f_\theta(s,a,z,s^+)\right].
    \label{eq:app-partition}
\end{equation}
Since $h_z^{\mathrm{tgt}}$ integrates to one, maximizing the negative cross-entropy in Equation~\eqref{eq:app-population-ce} is equivalent to
\begin{align}
    \mathcal J_{\mathrm{DFB}}(\theta;\bar h)
    ={}&
    (1-\gamma)\E_{\substack{z\sim p_Z\\(s,a,s')\sim\mathcal D}}\left[f_\theta(s,a,z,s')\right] \nonumber\\
    &+\gamma\E_{\substack{z\sim p_Z\\(s,a,s')\sim\mathcal D\\a'\sim\pi_z(\cdot\mid s')\\s^+\sim\rho}}\left[\bar h(s',a',z,s^+)f_\theta(s,a,z,s^+)\right] \nonumber\\
    &-\E_{\substack{z\sim p_Z\\(s,a)\sim\mathcal D_{SA}}}\left[\log Z_\theta(s,a,z)\right].
    \label{eq:app-score-objective}
\end{align}
In practice, $\bar h$ is represented by target networks and treated as stop-gradient during this update.


\subsection{Finite-pool approximation}
\label{app:finite-pool}

The population normalization in Equation~\eqref{eq:app-density-model} requires expectations under $\rho$. \name{} approximates these expectations using the next states in each minibatch.

Let
\[
    \{(s_i,a_i,s_i')\}_{i=1}^{n}\sim\mathcal D
\]
be a minibatch. The candidate states $\{s_j'\}_{j=1}^{n}$ are samples from the dataset next-state marginal $\rho$. For task index $z_i$, define
\begin{equation}
    f_{ij} := \frac{F_\theta(s_i,a_i,z_i)^\top B_\theta(s_j')}{\tau},
    \qquad
    p_{ij} := \frac{\exp f_{ij}}{\sum_{k=1}^{n}\exp f_{ik}}.
    \label{eq:app-finite-model}
\end{equation}
Thus $p_i$ is the model distribution over the empirical candidate pool.

Let
\begin{equation}
    a_i' := \sg\left[\pi_\omega(s_i',z_i)\right]
    \label{eq:app-bootstrap-action}
\end{equation}
and define the target-network distribution
\begin{equation}
    \bar p_{ij}' := \operatorname{softmax}_{j}\left(\frac{\bar F(s_i',a_i',z_i)^\top\bar B(s_j')}{\tau}\right).
    \label{eq:app-finite-target}
\end{equation}

The observed transition successor $s_i'$ provides the sample from the immediate transition distribution in Equation~\eqref{eq:app-normalized-bellman}. The finite-pool Bellman target is therefore
\begin{equation}
    p_{ij}^{\mathrm{tgt}} := \sg\left[(1-\gamma)\mathbf 1[j=i] + \gamma\bar p_{ij}'\right].
    \label{eq:app-finite-bellman}
\end{equation}
Both terms are normalized over $j$, so $\sum_jp_{ij}^{\mathrm{tgt}}=1$.

The empirical representation loss is
\begin{equation}
    \mathcal L_{\mathrm{occ}} := -\frac{1}{n}\sum_{i=1}^{n}\sum_{j=1}^{n}p_{ij}^{\mathrm{tgt}}\log p_{ij}.
    \label{eq:app-finite-loss}
\end{equation}
Equation~\eqref{eq:app-finite-loss} is the finite-support counterpart of the population cross-entropy in Equation~\eqref{eq:app-population-ce}. As the candidate pool grows, the batch softmax provides a self-normalized Monte Carlo approximation to the conditional density model in Equation~\eqref{eq:app-density-model}.


\subsection{Reward readout from the successor density}
\label{app:reward-readout}

The learned density also provides a value estimate for arbitrary state rewards. For a fixed policy $\pi_z$ and reward $r:\mathcal S\rightarrow\mathbb R$, define
\begin{equation}
    q_r^{\pi_z}(s,a) := \E_{s^+\sim\rho}\left[h_z^\pi(s,a,s^+)r(s^+)\right].
    \label{eq:app-true-reward-readout}
\end{equation}

\paragraph{Proposition 3 (Normalized reward value).}
The quantity in Equation~\eqref{eq:app-true-reward-readout} satisfies
\begin{equation}
    q_r^{\pi_z}(s,a) = (1-\gamma)Q_r^{\pi_z}(s,a),
    \label{eq:app-q-relation}
\end{equation}
where
\[
    Q_r^{\pi_z}(s,a) := \E_{\pi_z}\left[\sum_{t=0}^{\infty}\gamma^t r(s_{t+1}) \,\middle|\, s_0=s,\ a_0=a\right].
\]

\paragraph{Proof.}
By the definition of the density ratio,
\begin{align}
    q_r^{\pi_z}(s,a)
    &= \int_{\mathcal S}r(s^+)\,d\nu^{\pi_z}(s,a,s^+) \nonumber\\
    &= (1-\gamma)\int_{\mathcal S}r(s^+)\,dM^{\pi_z}(s,a,s^+) \nonumber\\
    &= (1-\gamma)\E_{\pi_z}\left[\sum_{t=0}^{\infty}\gamma^t r(s_{t+1}) \,\middle|\, s_0=s,\ a_0=a\right].
    \label{eq:app-reward-proof}
\end{align}
\hfill$\square$

\name{} replaces the true density ratio with the learned model:
\begin{equation}
    q_{\theta,r}(s,a,z) := \E_{s^+\sim\rho}\left[h_\theta(s,a,z,s^+)r(s^+)\right].
    \label{eq:app-model-reward-readout}
\end{equation}
If $h_\theta=h_z^\pi$, then $q_{\theta,r}=q_r^{\pi_z}=(1-\gamma)Q_r^{\pi_z}$. For an approximate successor model, $q_{\theta,r}$ is the corresponding model-based estimate.

Using a finite candidate pool $\{s_j^+\}_{j=1}^{P}$,
\begin{equation}
    q_{\theta,r}(s,a,z) \approx \sum_{j=1}^{P}p_j(s,a,z)r(s_j^+),
    \label{eq:app-finite-reward}
\end{equation}
where $p_j(s,a,z)$ is the normalized successor probability assigned by the model. Since $1-\gamma$ is independent of the action, exact density estimation preserves the action ranking induced by $Q_r^{\pi_z}$.


\section{Relation to Standard Forward-Backward Representations}
\label{app:fb-relation}

\subsection{Why the standard FB reward map does not transfer}
\label{app:fb-task-map}

Standard FB represents the density of the unnormalized successor measure using
\begin{equation}
    \frac{dM^{\pi_z}(s,a,\cdot)}{d\rho}(s^+) \approx F(s,a,z)^\top B(s^+).
    \label{eq:app-standard-fb}
\end{equation}
The linear factorization gives
\begin{align}
    \int_{\mathcal S}r(s^+)F(s,a,z)^\top B(s^+)\,d\rho(s^+)
    &= F(s,a,z)^\top\E_{s^+\sim\rho}\left[r(s^+)B(s^+)\right].
    \label{eq:app-standard-task}
\end{align}
This identity motivates the standard FB reward embedding
\[
    z_r \propto \E_\rho[rB].
\]

\name{} changes the role of the bilinear product. Under Equation~\eqref{eq:app-density-model}, $F^\top B/\tau$ is a log-score inside a normalized exponential, so the linear identity in Equation~\eqref{eq:app-standard-task} no longer holds.

There is also an exact ambiguity in the backward representation. Consider any constant vector $b\in\mathbb R^d$ and the transformation
\begin{equation}
    B(s^+) \mapsto B(s^+) + b.
    \label{eq:app-backward-shift}
\end{equation}
For fixed $(s,a,z)$, this changes every score by the same amount,
\[
    \frac{F(s,a,z)^\top(B(s^+)+b)}{\tau}
    =
    \frac{F(s,a,z)^\top B(s^+)}{\tau}
    +
    \frac{F(s,a,z)^\top b}{\tau}.
\]
The additive term does not depend on $s^+$ and therefore cancels under normalization:
\begin{equation}
    \frac{\exp(F^\top(B(s^+)+b)/\tau)}{\E_\rho[\exp(F^\top(B+b)/\tau)]}
    =
    \frac{\exp(F^\top B(s^+)/\tau)}{\E_\rho[\exp(F^\top B/\tau)]}.
    \label{eq:app-gauge-invariance}
\end{equation}

The standard FB reward moment is not invariant to the same transformation:
\begin{equation}
    \E_\rho\left[r(s^+)(B(s^+)+b)\right]
    =
    \E_\rho\left[r(s^+)B(s^+)\right]
    +
    b\,\E_\rho[r(s^+)].
    \label{eq:app-task-noninvariance}
\end{equation}
Hence two backward representations that define exactly the same \name{} successor distribution can produce different values of $\E_\rho[rB]$. The standard FB reward moment is therefore not an identifiable task index for the normalized model.

Centering the reward removes the particular shift term in Equation~\eqref{eq:app-task-noninvariance}, but it does not recover Equation~\eqref{eq:app-standard-task}: reward evaluation remains nonlinear in the backward features because $F^\top B$ appears inside a normalized exponential. \name{} therefore requires a task-selection mechanism distinct from the standard FB linear reward embedding.


\subsection{Reward-based selection within the learned policy family}
\label{app:policy-selection}

Rather than constructing a new task vector in backward-feature coordinates, \name{} keeps the task index in the same space used during pretraining. Let $\mathcal Z_{\mathrm{cand}}$ be a finite set of task indices drawn from the policy family covered during training. For a downstream reward $r$, define
\begin{equation}
    J_{\theta,r}(z) := \E_{s\sim\mu_{\mathrm{sel}}}\left[q_{\theta,r}(s,\pi_\omega(s,z),z)\right].
    \label{eq:app-policy-score}
\end{equation}
We select
\begin{equation}
    \hat z_r \in \arg\max_{z\in\mathcal Z_{\mathrm{cand}}}J_{\theta,r}(z)
    \label{eq:app-policy-selection}
\end{equation}
and execute the frozen policy $\pi_\omega(\cdot,\hat z_r)$.

This procedure uses the downstream reward only to evaluate policies through their predicted successor distributions. It therefore does not require the reward itself to share coordinates with the backward representation. The resulting guarantee is correspondingly weaker than the idealized standard FB interface: Equation~\eqref{eq:app-policy-selection} identifies the highest-valued policy within the covered candidate family, not necessarily a globally optimal policy for an arbitrary reward.